\section{Basic.lean}

\code{Basic.lean} is the repository foundation used by the V1 structured diff
layer. The structured diff modules import this foundation rather than
replacing it.

\subsection{Imports}

\begin{definition}[No imports]
\code{[Basic.lean]} has no imports. It uses Lean core declarations only.
\end{definition}

\subsection{Repository Primitives}

\begin{definition}[File]
\code{[File]} is an abbreviation for \code{String}. It denotes a file path or
file identifier.
\end{definition}

\begin{definition}[Content]
\code{[Content]} is an abbreviation for \code{String}. It denotes the complete
contents of a file as one string.
\end{definition}

\begin{definition}[Agent]
\code{[Agent]} is an abbreviation for \code{Nat}. It denotes an agent id.
\end{definition}

\begin{definition}[Repository]
\code{[Repository]} is an abbreviation for \code{File -> Option Content}. It is
a partial map from file names to contents.
\end{definition}

\begin{definition}[System]
\code{[System]} is a structure with fields \code{origin : Repository} and
\code{agents : Agent -> Repository}. It represents one central repository and
one local repository for each agent.
\end{definition}

\begin{definition}[Timeline]
\code{[Timeline]} is an inductive type indexed by \code{System}. It has two
constructors: \code{origin} and \code{agentic agent}.
\end{definition}

\subsection{Whole-File Patch Skeleton}

\begin{definition}[FilePatch]
\code{[FilePatch]} is an abbreviation for \code{File -> Option Content}. It is
the existing whole-file patch concept from the repository foundation.
\end{definition}

\begin{definition}[Commit]
\code{[Commit]} is a structure with fields \code{author : Agent},
\code{changes : FilePatch}, and \code{commitMessage : String}.
\end{definition}

\subsection{Existing Action Skeleton}

\code{Basic.lean} also contains an \code{ActionF}/\code{ActionM} skeleton for
repository-level actions. The structured diff layer leaves that skeleton in
place. It only changes the old whole-file patch name from \code{Diff} to
\code{FilePatch}, because the structured JSON diff introduced in
\code{Defns.lean} is named \code{Diff}.

\begin{definition}[Compatibility rename]
The compatibility edit is:
\[
  \code{Diff := File -> Option Content}
  \quad\longrightarrow\quad
  \code{FilePatch := File -> Option Content}.
\]
The \code{Commit.changes} field is updated accordingly.
\end{definition}

\subsection{Relationship To The Structured Diff Stack}

\begin{definition}[Separation of roles]
\code{Basic.lean} does not define structured line-edit declarations such as
\code{Range}, \code{Hunk}, or the JSON-facing \code{Diff}. Those belong to
\code{Defns.lean}. The older action skeleton in \code{Basic.lean} is preserved
as existing repository code and remains separate from the structured diff core.
\end{definition}
